% !TEX root = ../main.tex
\section{Conclusion}
\label{sec:conclusion}

We presented \medskillbench, the first skill-level execution benchmark for medical AI agents, covering 1,462 skills with eight-dimensional scoring in isolated sandbox environments.
Using this benchmark, we identified and quantified a systematic \emph{read-but-not-use} gap of 22 percentage points between skill invocation (96.2\%) and task completion (74.3\%).
Critically, stratification by \texttt{tool\_type} reveals that guide-style skills without executable interfaces exhibit the largest gap (0.269), while tool-based skills achieve substantially tighter coupling between comprehension and execution (0.069--0.111).
A two-stage decomposition of the \rbu gap into semantic grounding ($\delta_{\text{sem}}$) and execution completion ($\delta_{\text{exec}}$) components reveals that the dominant bottleneck is specification-to-action translation, not action-to-outcome execution.
This finding reframes the core bottleneck: the primary failure mode is not that models cannot operate tools, but that natural-language-only skill descriptions lack sufficient grounding for reliable execution.

Linking skills to real medical QA benchmarks, we further demonstrate that skill availability yields a measurable accuracy uplift of \todo{X} percentage points on average, with the effect strongly modulated by skill interface type---consistent with the \rbu gap analysis.
This closes the diagnostic loop: skill-level execution quality, as measured by \medskillbench, directly predicts downstream task performance gains.

Our results establish that medical agent skill quality can be systematically measured at the execution level, providing fine-grained diagnostic signal that task-level benchmarks cannot.
The evaluation-driven diagnostic paradigm extends naturally beyond the medical domain to any setting where LLM agents consume structured skill specifications, offering a principled path toward understanding and improving agent--tool interaction.
We discuss remaining limitations in \S\ref{sec:limitations}.
